% =====================================================================
% tikz-closure-flexibility-map.tex
% ---------------------------------------------------------------------
% Closure flexibility map: a 4 x 2 grid of constant-height text boxes
% (NOT a table). Columns = increasing flexibility (left to right);
% rows = level of aggregation (aggregate / sectoral).
%
%   row 1  aggregate : conventional CGE  fixed real wage  Leontief closure
%   row 2  sectoral  : no reallocation |  elastic labor supply models | [ -- ]
%                        (col 1)          (one box, cols 2-3)          (col 4)
%
% Double-headed arrows:
%   * "no reallocation" -> "conventional CGE": up first, then right
%     (one 90-degree corner); replaces the direct diagonal.
%   * the elastic box <-> its left neighbour "no reallocation" (short),
%     and <-> "fixed real wage" above it (short).
% Plain lines still link the aggregate row left to right.
%
% Drop-in: % =====================================================================
% tikz-closure-flexibility-map.tex
% ---------------------------------------------------------------------
% Closure flexibility map: a 4 x 2 grid of constant-height text boxes
% (NOT a table). Columns = increasing flexibility (left to right);
% rows = level of aggregation (aggregate / sectoral).
%
%   row 1  aggregate : conventional CGE  fixed real wage  Leontief closure
%   row 2  sectoral  : no reallocation |  elastic labor supply models | [ -- ]
%                        (col 1)          (one box, cols 2-3)          (col 4)
%
% Double-headed arrows:
%   * "no reallocation" -> "conventional CGE": up first, then right
%     (one 90-degree corner); replaces the direct diagonal.
%   * the elastic box <-> its left neighbour "no reallocation" (short),
%     and <-> "fixed real wage" above it (short).
% Plain lines still link the aggregate row left to right.
%
% Drop-in: % =====================================================================
% tikz-closure-flexibility-map.tex
% ---------------------------------------------------------------------
% Closure flexibility map: a 4 x 2 grid of constant-height text boxes
% (NOT a table). Columns = increasing flexibility (left to right);
% rows = level of aggregation (aggregate / sectoral).
%
%   row 1  aggregate : conventional CGE  fixed real wage  Leontief closure
%   row 2  sectoral  : no reallocation |  elastic labor supply models | [ -- ]
%                        (col 1)          (one box, cols 2-3)          (col 4)
%
% Double-headed arrows:
%   * "no reallocation" -> "conventional CGE": up first, then right
%     (one 90-degree corner); replaces the direct diagonal.
%   * the elastic box <-> its left neighbour "no reallocation" (short),
%     and <-> "fixed real wage" above it (short).
% Plain lines still link the aggregate row left to right.
%
% Drop-in: % =====================================================================
% tikz-closure-flexibility-map.tex
% ---------------------------------------------------------------------
% Closure flexibility map: a 4 x 2 grid of constant-height text boxes
% (NOT a table). Columns = increasing flexibility (left to right);
% rows = level of aggregation (aggregate / sectoral).
%
%   row 1  aggregate : conventional CGE  fixed real wage  Leontief closure
%   row 2  sectoral  : no reallocation |  elastic labor supply models | [ -- ]
%                        (col 1)          (one box, cols 2-3)          (col 4)
%
% Double-headed arrows:
%   * "no reallocation" -> "conventional CGE": up first, then right
%     (one 90-degree corner); replaces the direct diagonal.
%   * the elastic box <-> its left neighbour "no reallocation" (short),
%     and <-> "fixed real wage" above it (short).
% Plain lines still link the aggregate row left to right.
%
% Drop-in: \input{figures/tikz-closure-flexibility-map}
% Requires (already in the preamble): \usepackage{tikz} and
%   \usetikzlibrary{arrows.meta, positioning, calc}
% ---------------------------------------------------------------------
% Style = house TikZ idiom: sans-serif labels, rounded corners, light
% gray fills, Stealth arrowheads. The empty cell at row 2 / col 4 is a
% `placeholder`; the old row 1 / col 1 placeholder is gone (the elbow of
% the no-reallocation -> conventional-CGE arrow occupies that space).
% =====================================================================
\begin{figure}[htbp]
\centering
\begin{tikzpicture}[
    node distance = 1.4cm and 2.0cm,
    box/.style = {
      draw, rectangle, rounded corners=3pt,
      minimum width=3.4cm, minimum height=1.0cm,
      align=center, font=\sffamily\small,
      fill=gray!20, line width=0.5pt
    },
    widebox/.style = {
      box, minimum width=7.4cm          % spans two columns + one gap
    },
    placeholder/.style = {
      draw, rectangle, rounded corners=3pt,
      minimum width=3.4cm, minimum height=1.0cm,
      align=center, font=\sffamily\small,
      draw=gray!50, densely dashed, line width=0.4pt
    },
    rowlabel/.style = {font=\sffamily\small},
    axlabel/.style  = {font=\sffamily\small},
    link/.style     = {draw, thick},
    arr/.style      = {-{Stealth}, thick},
    flexarrow/.style = {{Stealth}-{Stealth}, thick}
  ]

  % ---- geometry -------------------------------------------------------
  \def\dx{4.0}    % horizontal pitch between column centres
  \def\rowy{1.9}  % vertical distance between row centres

  % ---- cells: row 1 = aggregate --------------------------------------
  \node[box] (r1c2) at (\dx,0)    {conventional CGE};
  \node[box] (r1c3) at (2*\dx,0)  {fixed real wage};
  \node[box] (r1c4) at (3*\dx,0)  {Leontief closure};

  % ---- cells: row 2 = sectoral ---------------------------------------
  \node[box]         (r2c1) at (0,-\rowy) {no reallocation};
  \node[widebox]     (ELS)  at (1.5*\dx,-\rowy) {elastic labor supply models};
  \node[placeholder] (r2c4) at (3*\dx,-\rowy) {};

  % ---- connector lines through the filled aggregate cells (left-right)
  \draw[link] (r1c2.east) -- (r1c3.west);
  \draw[link] (r1c3.east) -- (r1c4.west);

  % ---- double arrow: no reallocation -> conventional CGE
  %      up first, then right (one 90-degree corner)
  \draw[flexarrow] (r2c1.north) -- (0,0) -- (r1c2.west);

  % ---- short double arrows of the elastic-labour-supply box -----------
  \draw[flexarrow] (ELS.west)   -- (r2c1.east);          % to no reallocation
  \draw[flexarrow] (2*\dx,-1.4) -- (r1c3.south);         % up to fixed real wage

  % ---- outer axes -----------------------------------------------------
  % y-axis with the two row labels (rotated to avoid overlap)
  \draw[arr] (-2.9,-2.7) -- (-2.9,1.0);
  \draw (-2.9,0)        -- (-2.72,0);
  \draw (-2.9,-\rowy)   -- (-2.72,-\rowy);
  \node[rowlabel, rotate=90] at (-3.25,0)      {aggregate};
  \node[rowlabel, rotate=90] at (-3.25,-\rowy) {sectoral};

  % x-axis: flexibility, increasing to the right
  \draw[arr] (-2.9,-3.2) -- (14.4,-3.2);
  \node[axlabel, below] at (5.75,-3.25) {flexibility};

\end{tikzpicture}
\caption{A map of closure choices along the short-run (labour-market)
  axis. Columns order the closures by increasing flexibility; the
  aggregate row collects the single-wage closures, the sectoral row the
  sector-level ones. The double-width box holds the elastic-labour-supply
  family, tied by short double arrows to the no-reallocation and
  fixed-real-wage corners it interpolates between; the right-angled
  double arrow marks the reallocation margin from no reallocation up to
  the conventional CGE closure.}
\label{fig:flexibility-map}
\end{figure}
% Requires (already in the preamble): \usepackage{tikz} and
%   \usetikzlibrary{arrows.meta, positioning, calc}
% ---------------------------------------------------------------------
% Style = house TikZ idiom: sans-serif labels, rounded corners, light
% gray fills, Stealth arrowheads. The empty cell at row 2 / col 4 is a
% `placeholder`; the old row 1 / col 1 placeholder is gone (the elbow of
% the no-reallocation -> conventional-CGE arrow occupies that space).
% =====================================================================
\begin{figure}[htbp]
\centering
\begin{tikzpicture}[
    node distance = 1.4cm and 2.0cm,
    box/.style = {
      draw, rectangle, rounded corners=3pt,
      minimum width=3.4cm, minimum height=1.0cm,
      align=center, font=\sffamily\small,
      fill=gray!20, line width=0.5pt
    },
    widebox/.style = {
      box, minimum width=7.4cm          % spans two columns + one gap
    },
    placeholder/.style = {
      draw, rectangle, rounded corners=3pt,
      minimum width=3.4cm, minimum height=1.0cm,
      align=center, font=\sffamily\small,
      draw=gray!50, densely dashed, line width=0.4pt
    },
    rowlabel/.style = {font=\sffamily\small},
    axlabel/.style  = {font=\sffamily\small},
    link/.style     = {draw, thick},
    arr/.style      = {-{Stealth}, thick},
    flexarrow/.style = {{Stealth}-{Stealth}, thick}
  ]

  % ---- geometry -------------------------------------------------------
  \def\dx{4.0}    % horizontal pitch between column centres
  \def\rowy{1.9}  % vertical distance between row centres

  % ---- cells: row 1 = aggregate --------------------------------------
  \node[box] (r1c2) at (\dx,0)    {conventional CGE};
  \node[box] (r1c3) at (2*\dx,0)  {fixed real wage};
  \node[box] (r1c4) at (3*\dx,0)  {Leontief closure};

  % ---- cells: row 2 = sectoral ---------------------------------------
  \node[box]         (r2c1) at (0,-\rowy) {no reallocation};
  \node[widebox]     (ELS)  at (1.5*\dx,-\rowy) {elastic labor supply models};
  \node[placeholder] (r2c4) at (3*\dx,-\rowy) {};

  % ---- connector lines through the filled aggregate cells (left-right)
  \draw[link] (r1c2.east) -- (r1c3.west);
  \draw[link] (r1c3.east) -- (r1c4.west);

  % ---- double arrow: no reallocation -> conventional CGE
  %      up first, then right (one 90-degree corner)
  \draw[flexarrow] (r2c1.north) -- (0,0) -- (r1c2.west);

  % ---- short double arrows of the elastic-labour-supply box -----------
  \draw[flexarrow] (ELS.west)   -- (r2c1.east);          % to no reallocation
  \draw[flexarrow] (2*\dx,-1.4) -- (r1c3.south);         % up to fixed real wage

  % ---- outer axes -----------------------------------------------------
  % y-axis with the two row labels (rotated to avoid overlap)
  \draw[arr] (-2.9,-2.7) -- (-2.9,1.0);
  \draw (-2.9,0)        -- (-2.72,0);
  \draw (-2.9,-\rowy)   -- (-2.72,-\rowy);
  \node[rowlabel, rotate=90] at (-3.25,0)      {aggregate};
  \node[rowlabel, rotate=90] at (-3.25,-\rowy) {sectoral};

  % x-axis: flexibility, increasing to the right
  \draw[arr] (-2.9,-3.2) -- (14.4,-3.2);
  \node[axlabel, below] at (5.75,-3.25) {flexibility};

\end{tikzpicture}
\caption{A map of closure choices along the short-run (labour-market)
  axis. Columns order the closures by increasing flexibility; the
  aggregate row collects the single-wage closures, the sectoral row the
  sector-level ones. The double-width box holds the elastic-labour-supply
  family, tied by short double arrows to the no-reallocation and
  fixed-real-wage corners it interpolates between; the right-angled
  double arrow marks the reallocation margin from no reallocation up to
  the conventional CGE closure.}
\label{fig:flexibility-map}
\end{figure}
% Requires (already in the preamble): \usepackage{tikz} and
%   \usetikzlibrary{arrows.meta, positioning, calc}
% ---------------------------------------------------------------------
% Style = house TikZ idiom: sans-serif labels, rounded corners, light
% gray fills, Stealth arrowheads. The empty cell at row 2 / col 4 is a
% `placeholder`; the old row 1 / col 1 placeholder is gone (the elbow of
% the no-reallocation -> conventional-CGE arrow occupies that space).
% =====================================================================
\begin{figure}[htbp]
\centering
\begin{tikzpicture}[
    node distance = 1.4cm and 2.0cm,
    box/.style = {
      draw, rectangle, rounded corners=3pt,
      minimum width=3.4cm, minimum height=1.0cm,
      align=center, font=\sffamily\small,
      fill=gray!20, line width=0.5pt
    },
    widebox/.style = {
      box, minimum width=7.4cm          % spans two columns + one gap
    },
    placeholder/.style = {
      draw, rectangle, rounded corners=3pt,
      minimum width=3.4cm, minimum height=1.0cm,
      align=center, font=\sffamily\small,
      draw=gray!50, densely dashed, line width=0.4pt
    },
    rowlabel/.style = {font=\sffamily\small},
    axlabel/.style  = {font=\sffamily\small},
    link/.style     = {draw, thick},
    arr/.style      = {-{Stealth}, thick},
    flexarrow/.style = {{Stealth}-{Stealth}, thick}
  ]

  % ---- geometry -------------------------------------------------------
  \def\dx{4.0}    % horizontal pitch between column centres
  \def\rowy{1.9}  % vertical distance between row centres

  % ---- cells: row 1 = aggregate --------------------------------------
  \node[box] (r1c2) at (\dx,0)    {conventional CGE};
  \node[box] (r1c3) at (2*\dx,0)  {fixed real wage};
  \node[box] (r1c4) at (3*\dx,0)  {Leontief closure};

  % ---- cells: row 2 = sectoral ---------------------------------------
  \node[box]         (r2c1) at (0,-\rowy) {no reallocation};
  \node[widebox]     (ELS)  at (1.5*\dx,-\rowy) {elastic labor supply models};
  \node[placeholder] (r2c4) at (3*\dx,-\rowy) {};

  % ---- connector lines through the filled aggregate cells (left-right)
  \draw[link] (r1c2.east) -- (r1c3.west);
  \draw[link] (r1c3.east) -- (r1c4.west);

  % ---- double arrow: no reallocation -> conventional CGE
  %      up first, then right (one 90-degree corner)
  \draw[flexarrow] (r2c1.north) -- (0,0) -- (r1c2.west);

  % ---- short double arrows of the elastic-labour-supply box -----------
  \draw[flexarrow] (ELS.west)   -- (r2c1.east);          % to no reallocation
  \draw[flexarrow] (2*\dx,-1.4) -- (r1c3.south);         % up to fixed real wage

  % ---- outer axes -----------------------------------------------------
  % y-axis with the two row labels (rotated to avoid overlap)
  \draw[arr] (-2.9,-2.7) -- (-2.9,1.0);
  \draw (-2.9,0)        -- (-2.72,0);
  \draw (-2.9,-\rowy)   -- (-2.72,-\rowy);
  \node[rowlabel, rotate=90] at (-3.25,0)      {aggregate};
  \node[rowlabel, rotate=90] at (-3.25,-\rowy) {sectoral};

  % x-axis: flexibility, increasing to the right
  \draw[arr] (-2.9,-3.2) -- (14.4,-3.2);
  \node[axlabel, below] at (5.75,-3.25) {flexibility};

\end{tikzpicture}
\caption{A map of closure choices along the short-run (labour-market)
  axis. Columns order the closures by increasing flexibility; the
  aggregate row collects the single-wage closures, the sectoral row the
  sector-level ones. The double-width box holds the elastic-labour-supply
  family, tied by short double arrows to the no-reallocation and
  fixed-real-wage corners it interpolates between; the right-angled
  double arrow marks the reallocation margin from no reallocation up to
  the conventional CGE closure.}
\label{fig:flexibility-map}
\end{figure}
% Requires (already in the preamble): \usepackage{tikz} and
%   \usetikzlibrary{arrows.meta, positioning, calc}
% ---------------------------------------------------------------------
% Style = house TikZ idiom: sans-serif labels, rounded corners, light
% gray fills, Stealth arrowheads. The empty cell at row 2 / col 4 is a
% `placeholder`; the old row 1 / col 1 placeholder is gone (the elbow of
% the no-reallocation -> conventional-CGE arrow occupies that space).
% =====================================================================
\begin{figure}[htbp]
\centering
\begin{tikzpicture}[
    node distance = 1.4cm and 2.0cm,
    box/.style = {
      draw, rectangle, rounded corners=3pt,
      minimum width=3.4cm, minimum height=1.0cm,
      align=center, font=\sffamily\small,
      fill=gray!20, line width=0.5pt
    },
    widebox/.style = {
      box, minimum width=7.4cm          % spans two columns + one gap
    },
    placeholder/.style = {
      draw, rectangle, rounded corners=3pt,
      minimum width=3.4cm, minimum height=1.0cm,
      align=center, font=\sffamily\small,
      draw=gray!50, densely dashed, line width=0.4pt
    },
    rowlabel/.style = {font=\sffamily\small},
    axlabel/.style  = {font=\sffamily\small},
    link/.style     = {draw, thick},
    arr/.style      = {-{Stealth}, thick},
    flexarrow/.style = {{Stealth}-{Stealth}, thick}
  ]

  % ---- geometry -------------------------------------------------------
  \def\dx{4.0}    % horizontal pitch between column centres
  \def\rowy{1.9}  % vertical distance between row centres

  % ---- cells: row 1 = aggregate --------------------------------------
  \node[box] (r1c2) at (\dx,0)    {conventional CGE};
  \node[box] (r1c3) at (2*\dx,0)  {fixed real wage};
  \node[box] (r1c4) at (3*\dx,0)  {Leontief closure};

  % ---- cells: row 2 = sectoral ---------------------------------------
  \node[box]         (r2c1) at (0,-\rowy) {no reallocation};
  \node[widebox]     (ELS)  at (1.5*\dx,-\rowy) {elastic labor supply models};
  \node[placeholder] (r2c4) at (3*\dx,-\rowy) {};

  % ---- connector lines through the filled aggregate cells (left-right)
  \draw[link] (r1c2.east) -- (r1c3.west);
  \draw[link] (r1c3.east) -- (r1c4.west);

  % ---- double arrow: no reallocation -> conventional CGE
  %      up first, then right (one 90-degree corner)
  \draw[flexarrow] (r2c1.north) -- (0,0) -- (r1c2.west);

  % ---- short double arrows of the elastic-labour-supply box -----------
  \draw[flexarrow] (ELS.west)   -- (r2c1.east);          % to no reallocation
  \draw[flexarrow] (2*\dx,-1.4) -- (r1c3.south);         % up to fixed real wage

  % ---- outer axes -----------------------------------------------------
  % y-axis with the two row labels (rotated to avoid overlap)
  \draw[arr] (-2.9,-2.7) -- (-2.9,1.0);
  \draw (-2.9,0)        -- (-2.72,0);
  \draw (-2.9,-\rowy)   -- (-2.72,-\rowy);
  \node[rowlabel, rotate=90] at (-3.25,0)      {aggregate};
  \node[rowlabel, rotate=90] at (-3.25,-\rowy) {sectoral};

  % x-axis: flexibility, increasing to the right
  \draw[arr] (-2.9,-3.2) -- (14.4,-3.2);
  \node[axlabel, below] at (5.75,-3.25) {flexibility};

\end{tikzpicture}
\caption{A map of closure choices along the short-run (labour-market)
  axis. Columns order the closures by increasing flexibility; the
  aggregate row collects the single-wage closures, the sectoral row the
  sector-level ones. The double-width box holds the elastic-labour-supply
  family, tied by short double arrows to the no-reallocation and
  fixed-real-wage corners it interpolates between; the right-angled
  double arrow marks the reallocation margin from no reallocation up to
  the conventional CGE closure.}
\label{fig:flexibility-map}
\end{figure}